\documentclass[12pt,a4paper]{article}

\usepackage[spanish,es-tabla]{babel}
\usepackage[utf8]{inputenc}
\usepackage[T1]{fontenc}
\usepackage{lmodern}
\usepackage{geometry}
\usepackage{xcolor}
\usepackage{hyperref}
\usepackage{longtable}
\usepackage{array}
\usepackage{tikz}
\usepackage{enumitem}
\usetikzlibrary{arrows.meta,positioning,shapes.multipart,fit,backgrounds}

\geometry{margin=2.5cm}
\hypersetup{
    colorlinks=true,
    linkcolor=blue!60!black,
    urlcolor=blue!60!black
}

\title{\textbf{Implementacion de Patrones de Diseno en el Proyecto Bodas}\\
\large Tecnicas de Programacion}
\author{Victor Hernandez Sosa}
\date{28 de mayo de 2026}

\begin{document}

\maketitle
\tableofcontents
\newpage

\section{Introduccion}

El proyecto \textit{Bodas} es una aplicacion web desarrollada con Django para gestionar eventos, proveedores, pagos, subeventos, validaciones, reportes y notificaciones. La implementacion incorpora patrones de diseno para separar responsabilidades, mejorar la mantenibilidad y demostrar el uso de programacion orientada a objetos en un caso practico.

Los patrones no se han implementado solamente como ejemplos aislados. Algunos aparecen directamente en funcionalidades visibles para el usuario, como clonar eventos, contratar proveedores, generar reportes o validar eventos. Otros se usan como soporte interno para organizar la logica del sistema.

\section{Diagrama general}

\begin{figure}[h!]
\centering
\resizebox{\textwidth}{!}{
\begin{tikzpicture}[
    node distance=1.1cm and 1.6cm,
    box/.style={
        rectangle,
        rounded corners,
        draw=black!65,
        fill=blue!6,
        align=center,
        minimum width=3.2cm,
        minimum height=0.9cm
    },
    pattern/.style={
        rectangle,
        rounded corners,
        draw=black!60,
        fill=green!8,
        align=center,
        minimum width=3.2cm,
        minimum height=0.8cm
    },
    data/.style={
        rectangle,
        rounded corners,
        draw=black!60,
        fill=orange!10,
        align=center,
        minimum width=3.2cm,
        minimum height=0.8cm
    },
    arrow/.style={-{Latex[length=3mm]}, thick, draw=black!70}
]

\node[box] (user) {Usuario};
\node[box, below=of user] (templates) {Interfaz web\\Templates HTML};
\node[box, below=of templates] (views) {Vistas Django\\\texttt{view\_handlers.py}};
\node[data, below=of views] (models) {Modelos Django\\Evento, Proveedor, Pago};

\node[pattern, right=of views] (creation) {Creacionales\\Singleton, Factory,\\Abstract Factory, Builder, Prototype};
\node[pattern, below=of creation] (structural) {Estructurales\\Adapter, Decorator, Facade,\\Composite, Proxy, Bridge};
\node[pattern, below=of structural] (behavior) {Comportamiento\\Strategy, Observer, Command,\\State, Iterator, Template,\\Mediator, Chain, Memento, Visitor};

\node[data, left=of models] (db) {Base de datos\\SQLite / PostgreSQL};
\node[data, right=of models] (external) {Servicios externos\\Pagos, catering, streaming};

\draw[arrow] (user) -- (templates);
\draw[arrow] (templates) -- (views);
\draw[arrow] (views) -- (models);
\draw[arrow] (models) -- (db);
\draw[arrow] (views) -- (creation);
\draw[arrow] (views) -- (structural);
\draw[arrow] (views) -- (behavior);
\draw[arrow] (structural) -- (external);
\draw[arrow] (behavior) -- (models);
\draw[arrow] (creation) -- (models);

\end{tikzpicture}
}
\caption{Relacion general entre la interfaz, las vistas, los modelos y los patrones de diseno.}
\end{figure}

\section{Resumen de patrones implementados}

\begin{longtable}{|p{3.3cm}|p{4cm}|p{7cm}|}
\hline
\textbf{Patron} & \textbf{Modulo principal} & \textbf{Uso en el proyecto} \\
\hline
Singleton & \texttt{singleton.py} & Gestion centralizada de la configuracion global del sistema. \\
\hline
Factory Method & \texttt{factory\_method.py} & Creacion de proveedores concretos de pago, ubicacion o catering. \\
\hline
Abstract Factory & \texttt{abstract\_factory.py} & Creacion de familias de elementos relacionados segun el tipo de evento. \\
\hline
Builder & \texttt{builder.py} & Construccion guiada de eventos con diferentes configuraciones. \\
\hline
Prototype & \texttt{prototype.py} & Clonado de eventos existentes y sus subeventos. \\
\hline
Adapter & \texttt{adapter.py} & Unificacion de pasarelas de pago, catering y streaming. \\
\hline
Decorator & \texttt{decorator.py} & Adicion dinamica de extras al evento, como DJ, fotografo u open bar. \\
\hline
Facade & \texttt{facade.py} & Coordinacion de alta de evento, pago, proveedores y notificaciones. \\
\hline
Composite & \texttt{composite.py} & Representacion de eventos compuestos formados por subeventos. \\
\hline
Proxy & \texttt{proxy.py} & Control de acceso a operaciones segun el rol del usuario. \\
\hline
Bridge & \texttt{bridge.py} & Separacion entre tipo de reporte y formato de salida. \\
\hline
Strategy & \texttt{strategy.py} & Intercambio de algoritmos de filtrado y calculo. \\
\hline
Observer & \texttt{observer.py} & Notificacion de cambios en eventos, pagos y servicios. \\
\hline
Command & \texttt{command.py} & Encapsulacion de acciones sobre eventos como objetos ejecutables. \\
\hline
State & \texttt{state.py} & Gestion de estados del ciclo de vida de un evento. \\
\hline
Iterator & \texttt{iterator.py} & Recorrido normal, inverso, filtrado o paginado de colecciones. \\
\hline
Template Method & \texttt{template\_method.py} & Flujo comun de confirmacion con pasos especificos por tipo de evento. \\
\hline
Mediator & \texttt{mediator.py} & Comunicacion coordinada entre organizador, proveedor y coordinador. \\
\hline
Chain of Responsibility & \texttt{chain\_of\_responsibility.py} & Cadena de validadores para aprobar o rechazar eventos. \\
\hline
Memento & \texttt{memento.py} & Guardado y restauracion de estados de un evento. \\
\hline
Visitor & \texttt{visitor.py} & Operaciones externas sobre elementos de evento, como calculo o validacion. \\
\hline
\end{longtable}

\section{Explicacion detallada}

\subsection{Singleton}

El patron Singleton garantiza que exista una unica instancia de una clase. En el proyecto se aplica en \texttt{ConfiguracionGlobal}, que centraliza valores como moneda, impuestos, limite de asistentes, notificaciones activas y modo mantenimiento.

Se usa cuando una vista necesita consultar restricciones globales, por ejemplo al validar un evento. Su utilidad es evitar que distintas partes del sistema tengan configuraciones contradictorias.

\subsection{Factory Method}

Factory Method delega la creacion de objetos concretos a una fabrica. En el proyecto aparece en \texttt{ServiceFactory}, que crea proveedores de pago, ubicacion o catering segun un identificador.

Se usa cuando el sistema conoce el tipo de servicio solicitado, pero no quiere acoplar la vista o la logica principal a una clase concreta. Por ejemplo, pedir un proveedor \texttt{stripe} devuelve un objeto preparado para trabajar con Stripe.

\subsection{Abstract Factory}

Abstract Factory crea familias de objetos relacionados. En este proyecto permite generar componentes de un evento segun su categoria, como invitacion, decoracion y catering para bodas, conferencias o conciertos.

Se usa cuando todos los objetos deben pertenecer a la misma familia conceptual. Por ejemplo, una boda debe tener invitacion, decoracion y catering coherentes con una boda, no mezclados con los de una conferencia.

\subsection{Builder}

Builder permite construir objetos complejos paso a paso. En \texttt{builder.py} se usa para crear eventos con configuracion estandar, completa o personalizada.

Se usa en la creacion guiada de eventos. Es adecuado porque un evento tiene muchos campos y servicios opcionales: catering, escenario, iluminacion, seguridad, streaming y decoracion. El Builder evita constructores demasiado largos y mejora la legibilidad.

\subsection{Prototype}

Prototype crea nuevos objetos copiando uno existente. En el proyecto se usa para clonar eventos reales de la base de datos, incluyendo configuracion, servicios, decoradores y subeventos.

Se usa cuando el usuario quiere crear un evento similar a otro ya existente. En lugar de rellenar todos los datos desde cero, se clona el evento y se modifican los datos necesarios.

\subsection{Adapter}

Adapter convierte interfaces incompatibles en una interfaz comun. En \texttt{adapter.py} se adapta la comunicacion con pasarelas de pago, proveedores de catering y plataformas de streaming.

Se usa cuando el sistema debe tratar de forma uniforme servicios externos que funcionan de manera distinta. El usuario puede elegir Stripe, PayPal o MercadoPago, pero la aplicacion los procesa mediante una interfaz comun.

\subsection{Decorator}

Decorator permite anadir responsabilidades a un objeto sin modificar su clase original. En el proyecto se usa para agregar extras a eventos: musica en vivo, catering premium, seguridad VIP, streaming avanzado, DJ, iluminacion, open bar, valet parking, fotografo o videografo.

Se usa cuando el usuario activa o desactiva extras en el detalle de un evento. Cada decorador modifica la descripcion y aumenta el precio final.

\subsection{Facade}

Facade ofrece una interfaz simple para coordinar subsistemas complejos. En \texttt{facade.py}, \texttt{EventFacade} agrupa creacion de evento, pago, asignacion de proveedores y envio de notificaciones.

Se usa para representar operaciones completas de negocio sin exponer todos los pasos internos. Es util cuando una accion del sistema requiere coordinar varias clases.

\subsection{Composite}

Composite permite tratar objetos individuales y grupos de objetos de forma uniforme. En el proyecto representa eventos compuestos formados por subeventos.

Se usa cuando un evento principal contiene partes internas, por ejemplo ceremonia, banquete y fiesta posterior. El evento compuesto puede calcular presupuesto, duracion y capacidad total a partir de sus hijos.

\subsection{Proxy}

Proxy controla el acceso a un objeto real. En el proyecto se usa para simular distintos niveles de permiso: invitado, usuario, organizador y administrador.

Se usa cuando no todos los usuarios deben poder realizar las mismas acciones. Por ejemplo, un invitado puede ver informacion basica, pero un organizador puede editar o eliminar.

\subsection{Bridge}

Bridge separa una abstraccion de su implementacion. En el proyecto separa el tipo de reporte del formato de salida.

Se usa al generar reportes. El sistema puede producir un reporte resumen, detallado, financiero o completo, y combinarlo con formatos como JSON, CSV, HTML, PDF o email sin duplicar clases para cada combinacion.

\subsection{Strategy}

Strategy permite intercambiar algoritmos en tiempo de ejecucion. En el proyecto se aplica a filtros de eventos y calculos de coste.

Se usa cuando existen varias formas validas de resolver una misma tarea. Por ejemplo, filtrar por fecha, presupuesto, ubicacion o asistentes; o calcular costes con una estrategia estandar, premium o economica.

\subsection{Observer}

Observer permite que varios observadores reaccionen automaticamente cuando un objeto cambia. En el proyecto, cuando un evento cambia, se notifica a email, proveedores y analitica.

Se usa al actualizar presupuestos, contratar servicios, agregar extras o registrar pagos. La vista no necesita llamar manualmente a cada notificador; solo dispara el evento y los observadores reaccionan.

\subsection{Command}

Command encapsula una accion como un objeto. En el proyecto hay comandos para crear, actualizar, eliminar o aprobar eventos.

Se usa cuando se quiere guardar historial de acciones o permitir operaciones reversibles. Aunque en la interfaz no aparece como boton independiente, es una base adecuada para funcionalidades de auditoria o deshacer.

\subsection{State}

State permite que un objeto cambie su comportamiento segun su estado interno. En el proyecto modela estados como borrador, pendiente de aprobacion, aprobado, en progreso, completado o cancelado.

Se usa para controlar las transiciones validas de un evento. Por ejemplo, un evento en borrador puede pasar a pendiente de aprobacion, pero no deberia completarse directamente.

\subsection{Iterator}

Iterator proporciona una forma uniforme de recorrer colecciones sin exponer su estructura interna. En el proyecto se aplica a colecciones de eventos.

Se usa cuando se necesita recorrer eventos en orden normal, inverso, filtrado o paginado. Es especialmente util para separar la logica de recorrido de la logica de negocio.

\subsection{Template Method}

Template Method define un algoritmo general y permite que las subclases personalicen algunos pasos. En el proyecto se usa para confirmar eventos.

El flujo comun incluye validar datos, configurar servicios, calcular costes, ejecutar configuracion especifica y confirmar el evento. La parte especifica cambia segun el tipo: una conferencia necesita ponentes, una boda necesita tipo de ceremonia, un concierto necesita artistas, etc.

\subsection{Mediator}

Mediator centraliza la comunicacion entre objetos. En el proyecto coordina mensajes entre organizador, proveedor y coordinador.

Se usa para evitar que cada participante conozca directamente a todos los demas. En lugar de comunicar objetos punto a punto, todos hablan mediante el mediador.

\subsection{Chain of Responsibility}

Chain of Responsibility encadena varios validadores. Cada validador decide si aprueba la peticion o la rechaza, y si todo va bien pasa al siguiente.

En el proyecto se usa para validar eventos. La cadena comprueba capacidad, servicios requeridos, servicios externos, presupuesto, horarios y restricciones globales. Es visible en la pagina de validacion del evento.

\subsection{Memento}

Memento guarda el estado de un objeto para restaurarlo despues. En el proyecto se usa para crear puntos de control del estado de un evento.

Se usa cuando se quiere permitir restaurar una version anterior sin romper el encapsulamiento del objeto. Aunque no sea una funcion principal visible en la interfaz, es util para historiales, borradores o deshacer cambios.

\subsection{Visitor}

Visitor permite anadir operaciones sobre una jerarquia de objetos sin modificar sus clases. En el proyecto se usa para calcular costes, generar reportes o validar componentes como catering, streaming, decoracion y ubicacion.

Se usa cuando hay varios tipos de elementos y se quiere aplicar una operacion externa sobre todos ellos. Por ejemplo, calcular el coste total de diferentes componentes de un evento.

\section{Patrones visibles para el usuario}

Algunos patrones tienen una representacion clara en la interfaz:

\begin{itemize}[leftmargin=*]
    \item \textbf{Builder}: creacion guiada de eventos.
    \item \textbf{Prototype}: boton o flujo para clonar eventos.
    \item \textbf{Adapter}: contratacion de proveedores, pagos y streaming.
    \item \textbf{Decorator}: extras opcionales en el evento.
    \item \textbf{Composite}: gestion de subeventos.
    \item \textbf{Bridge}: descarga o generacion de reportes.
    \item \textbf{Chain of Responsibility}: validacion del evento.
    \item \textbf{Template Method}: confirmacion del evento segun su tipo.
\end{itemize}

\section{Patrones internos}

Otros patrones se utilizan principalmente para organizar el codigo:

\begin{itemize}[leftmargin=*]
    \item \textbf{Singleton}: configuracion global unica.
    \item \textbf{Factory Method}: creacion desacoplada de proveedores.
    \item \textbf{Abstract Factory}: familias coherentes de componentes.
    \item \textbf{Facade}: coordinacion de subsistemas.
    \item \textbf{Proxy}: control de acceso.
    \item \textbf{Strategy}: intercambio de algoritmos.
    \item \textbf{Observer}: notificaciones automaticas.
    \item \textbf{Command}: acciones encapsuladas.
    \item \textbf{State}: ciclo de vida del evento.
    \item \textbf{Iterator}: recorrido de colecciones.
    \item \textbf{Memento}: restauracion de estados.
    \item \textbf{Visitor}: operaciones externas sobre componentes.
    \item \textbf{Mediator}: comunicacion entre participantes.
\end{itemize}

\section{Conclusion}

La implementacion de patrones en el proyecto \textit{Bodas} permite demostrar el uso de programacion orientada a objetos en una aplicacion real. Los patrones creacionales ayudan a construir eventos y proveedores; los estructurales organizan la relacion con servicios externos, reportes y subeventos; y los de comportamiento gestionan validaciones, notificaciones, estados, flujos y comunicacion.

En conjunto, los patrones mejoran la separacion de responsabilidades, reducen el acoplamiento y hacen que el sistema sea mas facil de extender. Por ejemplo, se puede agregar una nueva pasarela de pago mediante un nuevo adaptador, un nuevo formato de reporte mediante Bridge, o una nueva validacion agregando otro elemento a la cadena de responsabilidad.

\end{document}
